\section{Results and Discussion}\label{sec:6}

This section reports benchmark results on the demo corpus defined in Section~\ref{sec:5}, compares ranking methods, and discusses the limits of the evidence.
All figures come from committed benchmark artifacts; they do not predict hiring outcomes at production scale.

\subsection{Evaluation scope and protocol}\label{sec:6-evaluated}

We evaluated \emph{resume-to-job matching}: for each of 30 sample resumes, the system ranks 15 job postings and is judged at cutoff $K{=}5$.
The 47 manually labeled pairs use the relevance scale defined in Section~\ref{sec:5-corpus}.
Metrics are macro-averaged across all 30 resumes so that no single profile dominates the score.

\subsection{Baseline methods}\label{sec:6-baselines}\label{sec:6-comparison}

The first comparison contrasts lexical baselines, dense semantic matching, skill overlap, hybrid semantic--skill matching, and the portal-default composite scorer on the same labeled set.
Multimodal weighted scoring achieves the highest ranking quality (nDCG@5 $= 0.924$) among the methods listed, while BM25 and TF--IDF remain close (nDCG@5 $\approx 0.89$--$0.90$).
On this small job pool, lexical baselines remain competitive, so gains from dense and hybrid matching are measurable but modest.
Semantic similarity alone (nDCG@5 $= 0.878$) underperforms both keyword baselines and multimodal blends, supporting the use of explicit skills alongside embeddings.
The full method-comparison table is provided in the supplementary information.

\subsection{Main results}\label{sec:6-main-result}\label{sec:6-progression}\label{sec:6-fusion}

The progression study (Table~\ref{tab:progression}) ranks every job for every resume and tracks how design choices accumulate on the same labeled set.
The portal-default six-channel composite reaches nDCG@5 $= 0.949$; the strongest \emph{single} configuration (multimodal weighted scoring) reaches $0.924$, and plain semantic cosine reaches $0.878$.
Reciprocal-rank fusion of four lists ($0.913$) helps diversity but does not beat the composite on this corpus.
Cross-encoder reranking improves over bare semantic search ($0.939$) yet still does not exceed the composite while adding roughly $340\times$ the compute, so it remains disabled by default.
Because the job pool is only 15 and the labels are exclusively positive, recall@5 saturates (many methods $\geq 0.93$), so the nDCG@5 differences are informative but narrow.

\begin{center}
\refstepcounter{table}\label{tab:progression}
\par\smallskip\noindent\textbf{Table~\thetable:} Paper progression benchmark (regression baselines, $K{=}5$).
\par\smallskip
\begin{tabular}{@{}lrrr@{}}
\toprule
Method & P@5 & R@5 & nDCG@5 \\
\midrule
Production composite (6-channel, default) & 0.293 & 0.933 & 0.949 \\
Cross-encoder reranker (disabled by default) & 0.307 & 0.983 & 0.939 \\
Multimodal weighted blend & 0.287 & 0.933 & 0.924 \\
RRF ensemble (4 lists) & 0.293 & 0.950 & 0.913 \\
TF-IDF cosine (lexical) & 0.307 & 0.983 & 0.905 \\
BM25 (lexical) & 0.307 & 0.983 & 0.902 \\
Semantic cosine & 0.267 & 0.867 & 0.878 \\
Soft skill embedding & 0.280 & 0.900 & 0.869 \\
Skills Jaccard & 0.233 & 0.733 & 0.748 \\
\bottomrule
\end{tabular}
\par\vspace{0.4em}
\footnotesize{Source: committed benchmark artifacts \path{comparison_table.json} (method variants) and \path{composite_eval_report.json} (portal composite); resume-to-jobs, $K{=}5$, 30 resumes $\times$ 15 jobs, gated by the $\pm 0.04$ nDCG regression tolerance. All pairwise nDCG@5 differences lie within the small-corpus confidence intervals --- no method is statistically distinguishable (see \emph{Statistical significance}).}
\end{center}

Learned logistic-regression fusion, evaluated held-out, reaches nDCG@5 $= 0.917$ and does not beat the fixed composite; its in-sample fit is higher but, because it is fit and judged on the same labeled pairs, we report only the held-out value and treat the in-sample number as an optimistic upper bound rather than a result. A protocol-gated search over 25 configurations finds no configuration that significantly beats the fixed composite after Holm correction, so the auditable fixed composite is retained.
Applying hard constraints after scoring lowers nDCG@5 to $0.908$ because filters can remove jobs that annotators marked as partially relevant.
Full fusion and model-parameter tables are provided in the supplementary information.

\leadpara{Statistical significance}
On this small corpus ($n{=}30$ resumes) the ranking differences are \emph{not} statistically robust. The composite's improvement over semantic-only ranking is positive but not significant (paired bootstrap $N{=}5000$, seed~42: $\Delta{=}{+}0.071$, two-sided $p{=}0.10$, 95\% CI crossing zero), and no pairwise comparison survives Holm--Bonferroni correction across the method family. We therefore report ranking \emph{parity} --- no method is statistically distinguishable from the others on this corpus --- rather than superiority; the lexical baselines (BM25, TF--IDF) are likewise not significantly different from semantic cosine. This parity is partly an instrument limitation: the labels are exclusively positive and the 15-job pool is barely larger than the cutoff $K{=}5$, so nDCG@5 has limited power to separate methods (a lexical baseline already reaches 0.905). We do not over-interpret parity as method equivalence; a larger, explicitly-negative-judged benchmark is required to discriminate ranking methods.

The main finding is modest but consistent: matching improves when the system combines language semantics with explicit skills, and the highest scores arise from hybrid signals rather than either channel alone.
No method achieves high precision---at best, roughly three in ten top-five slots are relevant (P@5 $\approx 0.31$)---because the corpus is small, labels are sparse, and several jobs can appear plausible for the same resume.
We therefore characterize JobMatch as \emph{assistive} ranking with explainable reasons, not as an automated hiring decision system.

\subsection{Multi-agent state sharing and usability}\label{sec:6-agent-sharing}\label{sec:6-agentic}

When a candidate edits a resume or an employer revises a job description, the owning agent rebuilds a versioned snapshot and shared embedding, then notifies the Matchmaking layer through the event bus (Section~\ref{sec:3}--Section~\ref{sec:4}).
The event-driven path addresses a usability problem seen in monolithic ATS tools: rankings that stay stale after a profile change, or that change without a visible cause.
Role-separated agents also keep each side's data under its owner's control, which supports trust and auditability even when two users see the same match score from the neutral Matchmaking agent.
The benchmarks validate that the shared scoring core ranks sensibly under manual labels; they do not yet prove that multi-agent coordination reduces time-to-hire or increases applicant satisfaction in the field.

\subsection{Strengths}\label{sec:6-strengths}\label{sec:6-ablation-latency}

Any method that combines text meaning with skills beats semantic-only search (nDCG@5 $0.878$ vs.\ $0.917$--$0.949$ for richer composites).
The portal-default composite reaches nDCG@5 $=0.949$; semantic+skills alone reaches $0.917$, while using experience, compensation, or location alone yields poor rankings (nDCG@5 $\leq 0.39$).
Detailed ablation rows are provided in the supplementary information.

Lightweight scorers complete in well under one millisecond per resume on 15 jobs; the full composite bi-encoder run averages ${\sim}0.4$\,ms per query.
That headroom matters for interactive portals where a user expects refreshed matches after saving a profile.
Cross-encoder reranking is much slower (${\sim}141.7$\,ms/query) and lowers nDCG@5 in the committed report, so it remains disabled by default.
The full latency table is provided in the supplementary information.

Hard-negative mining finds zero overlap between declared relevant jobs and 150 high-scoring mined negatives, supporting internal label consistency.
The live portal defaults to composite scoring because each channel maps to a human-readable explanation, even though some research variants score slightly higher offline.

\subsection{Limitations}\label{sec:6-limitations}\label{sec:6-fairness}\label{sec:6-hard-negs}\label{sec:6-latency}

The evaluation has four main limits: the labeled corpus is small, the fairness audit is synthetic, no live-user study was run, and JobMatch remains a research prototype rather than a validated deployment.
Thirty resumes and fifteen jobs make recall@5 easy to saturate---many methods reach ${\ge}0.93$---so differences in nDCG are informative but narrow.
Results may not transfer to employers with thousands of open roles or to domains unlike our demo profiles.
Even the best methods place relevant jobs in only about three of every ten top-five slots (P@5 $\approx 0.31$).
Users will still see imperfect suggestions and must rely on explanations and their own judgment.

Cross-encoder reranking is ${\sim}340\times$ slower than composite bi-encoder scoring and does not help on this corpus: as a standalone ranker it reaches nDCG@5 $= 0.939$ (below the composite's $0.949$), and using it to rerank the composite's top-20 pool \emph{lowers} nDCG@5 to $0.834$ ($\Delta{=}{-}0.108$)---consistent with leaving it disabled by default.
Learned fusion and calibration fit on the same labeled pairs they are judged against, which risks overfitting.

The fairness audit is narrow: ten fabricated profile pairs undergo controlled name, summary, or metadata edits.
Seven pairs trigger flags, all with the top-1 job stable, mostly due to small rank shifts or score changes; the maximum score shift is $0.017$.
A complementary proxy-group check (\path{fairness_eval.py}) reports selection-rate disparate-impact ratios of $0.82$ across experience tiers and $0.75$ across remote preference (1.0 is parity), committed to \path{benchmark_outputs/fairness_eval.json}.
These checks are useful engineering probes, but they do not certify demographic fairness or substitute for audits on larger labeled populations.

We did not measure interview rates, time-to-fill, applicant satisfaction, recruiter workload, or live-user behavior in the portals.
The implementation should therefore be read as a research prototype with reproducible offline evidence, not as a validated deployment study or production hiring system.
Both explainers are perfectly consistent across synthetic profile variants and never make unsupported component claims, but they cite an explicit matched or missing skill in only about a quarter of bullets, so most explanations are flagged for low specificity---motivating richer explanation templates.
Detailed fairness and explainability audit tables are provided in the supplementary information.

\subsection{Implications for users}\label{sec:6-implications}

JobMatch is designed as an assistive representative: maintain a structured profile, inspect why a job ranked highly or poorly, and decide whether to save or apply.
The evaluation indicates that suggestions are most reliable when the system weights both resume text and listed skills---consistent with how applicants skim postings for fit.
Top-five lists are neither exhaustive nor authoritative; with P@5 near 0.31, several shown jobs may be weak fits, and relevant openings can be missed when the job pool exceeds fifteen roles.

\leadpara{Employers}
The employer portal provides symmetric reverse matching: ranked candidates with component-level reasons instead of an opaque score.
Hybrid ranking supports recruiters who weight keywords and hiring managers who read career narrative.
Hard post-filters can remove partially suitable candidates from the graded top five, so the portal keeps constraints optional and off by default.
When a job description changes, the Employer Agent update path refreshes rankings without re-upload to a separate scoring service.

Both sides benefit when software explains its suggestions and leaves hire, reject, and apply decisions with users.
The benchmarks support that design on a demo corpus: combined signals rank better than single-signal search, explanations align with composite channels, and the multi-agent layout keeps profile edits on the owning side while a neutral Matchmaking Agent serves both portals.
They do not establish production gains on hiring outcomes; that requires larger labeled datasets, portal-default ablations, and field studies.
